\chapter{Diseases of the Throat and Oral Cavity}
\label{ch:throat}
The throat (pharynx) and oral cavity serve essential functions in respiration, deglutition, speech, and taste. The oral cavity includes the lips, tongue, gingiva, floor of mouth, palate, and tonsils. The pharynx connects the nasal and oral cavities to the larynx and esophagus. Diseases range from common infections to malignant neoplasms.

%-------------------------------------------------------------
\section{Infectious Diseases}
%-------------------------------------------------------------

%-------------------------------------------------------------

%-------------------------------------------------------------

%-------------------------------------------------------------

%-------------------------------------------------------------

%-------------------------------------------------------------

%-------------------------------------------------------------

\label{sec:Infectious Diseases}
%-------------------------------------------------------------%-------------------------------------------------------------

\subsection{Diphtheria}
\label{sec:throat:Diphtheria}
Diphtheria is a serious bacterial infection caused by \textit{Corynebacterium diphtheriae} that produces an exotoxin causing tissue necrosis and systemic toxicity. The condition results from infection with toxigenic strains of \textit{C. diphtheriae}. Clinical features include a tough, gray-white pseudomembrane on the pharynx and tonsils, neck swelling (``bull neck''), fever, and toxic appearance; the exotoxin can cause myocarditis and neuropathy. Diagnosis is by culture on tellurite agar and toxin detection. Management includes diphtheria antitoxin and antibiotics (penicillin or erythromycin). Prognosis has been dramatically improved by vaccination with diphtheria toxoid (as part of DTaP/Tdap), which has made the disease rare in developed countries.

\subsection{Laryngitis}
\label{sec:Laryngitis}
Laryngitis is an inflammation of the larynx that presents with hoarseness (dysphonia), most commonly viral in acute cases. The condition results from vocal cord inflammation due to viral infection, voice overuse (vocal cord nodules), gastroesophageal reflux disease (laryngopharyngeal reflux), smoking, allergies, or chronic infections. Clinical features include hoarseness, sore throat, cough, and sometimes low-grade fever; chronic laryngitis lasts more than 3 weeks. Diagnosis is by laryngoscopy showing vocal cord erythema and edema. Management includes voice rest, hydration, humidification, and addressing underlying causes; persistent hoarseness requires laryngoscopy to exclude vocal cord polyps, cysts, or malignancy. Prognosis is excellent for acute laryngitis, with resolution within 1--2 weeks.

\subsection{Pharyngitis}
\label{sec:Pharyngitis}
Pharyngitis is an inflammation of the pharynx, most commonly caused by viral infections (rhinoviruses, coronaviruses, adenovirus, influenza, Epstein-Barr virus); group A \textit{Streptococcus pyogenes} (GAS) is the most important bacterial cause, accounting for 15--30\% of pharyngitis cases in children and 5--15\% in adults. The condition results from infectious inflammation of the pharyngeal mucosa. Clinical features of GAS pharyngitis include acute onset of sore throat, fever, tonsillar exudates, tender anterior cervical lymphadenopathy, and palatal petechiae, typically without cough. Diagnosis is based on the Centor/McIsaac score, rapid antigen detection testing, and throat culture. Management involves penicillin or amoxicillin for GAS pharyngitis to prevent acute rheumatic fever. Prognosis is excellent with appropriate treatment; complications include peritonsillar abscess and poststreptococcal glomerulonephritis.

\subsection{Tonsillitis}
\label{sec:Tonsillitis}
Tonsillitis is an inflammation of the palatine tonsils, which are lymphoid tissues forming part of Waldeyer's ring, and may be acute (viral or bacterial, with GAS being the most common bacterial cause) or chronic (persistent low-grade infection with recurrent exacerbations). The condition results from infectious inflammation of the tonsillar tissue. Clinical features include sore throat, odynophagia, fever, halitosis, and cervical lymphadenopathy; tonsillar hypertrophy may cause obstructive sleep-disordered breathing. Diagnosis is based on clinical examination. Management of recurrent tonsillitis follows the Paradise criteria for tonsillectomy; peritonsillar abscess (quinsy) requires needle aspiration or incision and drainage with antibiotic therapy. Prognosis is excellent with appropriate treatment.

%-------------------------------------------------------------
\section{Obstructive and Sleep-Related Disorders}
%-------------------------------------------------------------

%-------------------------------------------------------------

%-------------------------------------------------------------

%-------------------------------------------------------------

%-------------------------------------------------------------

%-------------------------------------------------------------

%-------------------------------------------------------------

\label{sec:Obstructive and Sleep-Related Disorders}
%-------------------------------------------------------------%-------------------------------------------------------------

\subsection{Obstructive Sleep Apnea}
\label{sec:throat:Obstructive Sleep Apnea}
Obstructive sleep apnea is a disorder characterized by repetitive episodes of partial or complete upper airway obstruction during sleep, resulting in apneas and hypopneas. Risk factors include obesity, male sex, large neck circumference, retrognathia, tonsillar and adenoid hypertrophy, and craniofacial abnormalities. The condition results from collapse of the pharyngeal airway during sleep. Clinical features include loud snoring, witnessed apneas, excessive daytime sleepiness, morning headaches, and cognitive impairment; untreated OSA increases the risk of hypertension, atrial fibrillation, stroke, heart failure, and metabolic syndrome. Diagnosis is by polysomnography measuring the apnea-hypopnea index (AHI). Management includes continuous positive airway pressure (CPAP), weight loss, oral appliances, and surgical options (uvulopalatopharyngoplasty, maxillomandibular advancement, hypoglossal nerve stimulation). Prognosis is favorable with effective treatment and adherence.

\subsection{Tonsillar Hypertrophy}
\label{sec:Tonsillar Hypertrophy}
Tonsillar hypertrophy is an enlargement of the palatine and adenoid tonsils that is common in children and is the most common cause of obstructive sleep apnea in the pediatric population. The condition results from lymphoid hyperplasia, often in response to recurrent infections. Clinical features include mouth breathing, snoring, restless sleep, recurrent otitis media, and difficulty swallowing; severe hypertrophy may cause failure to thrive and cor pulmonale. Diagnosis is based on clinical examination and history. Management involves tonsillectomy with or without adenoidectomy. Prognosis is excellent following surgical intervention.

%-------------------------------------------------------------
\section{Oral Cavity Diseases}
%-------------------------------------------------------------

%-------------------------------------------------------------

%-------------------------------------------------------------

%-------------------------------------------------------------

%-------------------------------------------------------------

%-------------------------------------------------------------

%-------------------------------------------------------------

\label{sec:Oral Cavity Diseases}
%-------------------------------------------------------------%-------------------------------------------------------------

\subsection{Aphthous Ulcers (Canker Sores)}
\label{sec:Aphthous Ulcers (Canker Sores)}
Aphthous ulcers are recurrent, painful, shallow ulcers of the oral mucosa that affect up to 25\% of the population; they are not contagious and are distinct from cold sores (herpes labialis). The condition results from a multifactorial etiology including local trauma, stress, nutritional deficiencies (iron, folate, B12), hormonal changes, and immune dysregulation. Clinical forms include minor aphthous (most common, small ulcers less than 1 cm, heal without scarring), major aphthous (larger, deeper ulcers, may scar), and herpetiform aphthous (multiple pinpoint ulcers that may coalesce). Diagnosis is clinical. Management is symptomatic, including topical anesthetics (lidocaine), antimicrobial rinses (chlorhexidine), topical corticosteroids (triamcinolone paste), and for severe cases, systemic agents (colchicine, thalidomide, apremilast). Prognosis is good, with lesions healing spontaneously but frequently recurring.

\subsection{Gingivitis and Periodontitis}
\label{sec:Gingivitis and Periodontitis}
Gingivitis is an inflammation of the gingiva without attachment loss, representing the earliest and most common form of periodontal disease, caused by dental plaque biofilm accumulation. Periodontitis extends beyond the gingiva with loss of connective tissue attachment and alveolar bone, leading to pocket formation, gingival recession, tooth mobility, and eventual tooth loss. The condition results from bacterial plaque accumulation and the host inflammatory response. Risk factors include smoking, diabetes mellitus, genetic susceptibility, and immunodeficiency. Clinical features include red, swollen, bleeding gingiva for gingivitis; periodontitis additionally presents with pocket formation, gingival recession, and tooth mobility. Diagnosis is based on clinical examination and probing depths. Management of gingivitis involves oral hygiene instruction and professional cleaning; periodontitis treatment involves scaling and root planing, surgical intervention, and systemic antibiotics for aggressive disease. Prognosis is excellent for gingivitis with proper hygiene; periodontitis requires long-term maintenance.

\subsection{Oral Candidiasis}
\label{sec:Oral Candidiasis}
Oral candidiasis is a fungal infection of the oral mucosa caused by \textit{Candida} species, most commonly \textit{Candida albicans}, with predisposing factors including immunosuppression (HIV/AIDS, chemotherapy), diabetes, denture use, inhaled corticosteroids, xerostomia, and broad-spectrum antibiotics. The condition results from overgrowth of \textit{Candida} organisms in the oral cavity. Clinical features include acute pseudomembranous candidiasis presenting with white, removable plaques on the buccal mucosa, tongue, palate, and oropharynx, revealing an erythematous base when scraped; erythematous (atrophic) and chronic hyperplastic forms also exist. Diagnosis is confirmed by wet mount with potassium hydroxide showing pseudohyphae and blastospores. Management involves topical antifungals (nystatin, clotrimazole) for mild disease or systemic azoles (fluconazole) for moderate-to-severe or refractory disease. Prognosis is excellent with appropriate antifungal therapy.

\subsection{Oral Leukoplakia}
\label{sec:Oral Leukoplakia}
Oral leukoplakia is a white plaque of the oral mucosa that cannot be scraped off and does not correspond to any other diagnosable condition; it is the most common oral premalignant lesion, with a malignant transformation rate of approximately 3--5\%. Risk factors include tobacco use (smoking and smokeless), alcohol consumption, and betel nut chewing. The condition results from chronic mucosal irritation leading to epithelial dysplasia. Homogeneous leukoplakia has a lower malignant potential than non-homogeneous (speckled) leukoplakia. Clinical features include a white plaque on the oral mucosa. Diagnosis requires biopsy to assess for epithelial dysplasia. Management depends on the degree of dysplasia: observation for mild dysplasia, surgical excision for moderate-to-severe dysplasia, and close follow-up for all patients with cessation of risk factors. Prognosis is good with early detection and appropriate management.

\subsection{Oral Lichen Planus}
\label{sec:Oral Lichen Planus}
Oral lichen planus is a chronic inflammatory condition of the oral mucosa mediated by T-cell autoimmune attack on basal keratinocytes, affecting approximately 1--2\% of the population, most commonly middle-aged women. The condition results from autoimmune T-cell-mediated damage to basal keratinocytes. Clinical features include the reticular form with characteristic Wickham striae (white, lacy patches) on the buccal mucosa, tongue, and gingiva, usually asymptomatic; the erosive form causes painful, erythematous ulcerated areas. There is a small (approximately 1\%) risk of malignant transformation. Diagnosis is by biopsy showing a band-like lymphocytic infiltrate at the epithelial-connective tissue interface with basal cell degeneration. Management for symptomatic disease includes topical corticosteroids and calcineurin inhibitors; long-term surveillance is recommended. Prognosis is generally good, but the condition is chronic and requires surveillance for malignant transformation.

\subsection{Oral Squamous Cell Carcinoma}
\label{sec:Oral Squamous Cell Carcinoma}
Oral squamous cell carcinoma is a malignant neoplasm that accounts for over 90\% of oral malignancies, with approximately 370,000 new cases annually worldwide, most commonly arising on the lateral tongue and floor of the mouth. The condition results from malignant transformation of oral epithelial cells, driven by tobacco use, heavy alcohol consumption, betel nut chewing, and HPV infection. Clinical features include a non-healing ulcer, indurated mass, or white or red plaque that may be painless initially; advanced disease causes pain, bleeding, dysphagia, and trismus. Diagnosis is by incisional biopsy with TNM staging. Management of early disease is surgical excision with or without neck dissection; advanced disease requires multimodal therapy with surgery, radiation, and chemotherapy. Prognosis is related to stage at diagnosis, with 5-year survival rates of approximately 80--90\% for stage I and less than 30\% for stage IV.

%-------------------------------------------------------------
\section{Voice Disorders}
%-------------------------------------------------------------

%-------------------------------------------------------------

%-------------------------------------------------------------

Voice disorders, or dysphonia, involve abnormal alterations in vocal pitch, loudness, quality, or vocal effort resulting from laryngeal pathology. Etiologies encompass organic vocal fold lesions (nodules, polyps, granulomas), neurological impairment (vocal fold paralysis), and functional muscle tension dysphonia. Multidisciplinary management involves videostroboscopic evaluation, targeted voice therapy, medical optimization, and micro-laryngeal surgery.

%-------------------------------------------------------------

%-------------------------------------------------------------

%-------------------------------------------------------------

%-------------------------------------------------------------

\label{sec:Voice Disorders}
%-------------------------------------------------------------%-------------------------------------------------------------

\subsection{Vocal Cord Nodules}
\label{sec:Vocal Cord Nodules}
Vocal cord nodules are bilateral, symmetric, callus-like thickenings of the true vocal folds at the junction of the anterior and middle thirds, the point of maximal vibratory impact, and are the most common cause of chronic hoarseness in adults. The condition results from chronic vocal abuse leading to repetitive trauma of the vocal folds. Clinical features include hoarseness, breathy voice, vocal fatigue, and decreased vocal range. Diagnosis is by laryngoscopy or videostroboscopy showing smooth, bilateral masses. Management involves voice therapy as conservative management, with surgical excision reserved for refractory cases. Prognosis is excellent with appropriate voice therapy.

\subsection{Vocal Cord Paralysis}
\label{sec:Vocal Cord Paralysis}
Vocal cord paralysis is a condition resulting from disruption of the recurrent laryngeal nerve, superior laryngeal nerve, or vagus nerve, with the most common cause being iatrogenic (thyroid surgery, thoracic and neck surgery); other causes include tumors, stroke, viral infections, and idiopathic etiology. The condition results from nerve injury leading to loss of vocal fold movement. Clinical features of unilateral paralysis include hoarseness, aspiration, and a breathy voice; bilateral paralysis presents with stridor and respiratory distress. Diagnosis is by flexible laryngoscopy showing an immobile vocal cord. Management includes voice therapy, injection laryngoplasty, medialization thyroplasty for unilateral paralysis, and tracheostomy or posterior cordotomy for bilateral paralysis with airway compromise. Prognosis depends on the underlying cause and severity.

